\subsection{Error Analysis}
\label{subsec:error-analysis}

To complement the aggregate evaluation metrics, a detailed error analysis was
performed on the held-out test set for all four models and all five spectral
modalities. The analysis focuses on the missed-positive error rate for each
functional group, defined as the proportion of false negatives relative to the
number of positive samples for that label. This perspective is particularly
informative in the present multi-label setting, because global metrics such as
micro-F1 can remain comparatively high even when individual functional groups
are predicted poorly. Reporting the error analysis on the test set is
methodologically preferable to using the validation set, as the latter was used
during model selection and threshold tuning and therefore does not provide the
same degree of independence.

\begin{figure}[t]
    \centering
    \begin{minipage}[t]{0.49\linewidth}
        \centering
        \includegraphics[width=\linewidth]{images/Error_Analysis/CNN_IR_test_error_analysis_detailed.png}
    \end{minipage}\hfill
    \begin{minipage}[t]{0.49\linewidth}
        \centering
        \includegraphics[width=\linewidth]{images/Error_Analysis/MPS_IR_test_error_analysis_detailed.png}
    \end{minipage}

    \vspace{0.5em}

    \begin{minipage}[t]{0.49\linewidth}
        \centering
        \includegraphics[width=\linewidth]{images/Error_Analysis/TTN_IR_test_error_analysis_detailed.png}
    \end{minipage}\hfill
    \begin{minipage}[t]{0.49\linewidth}
        \centering
        \includegraphics[width=\linewidth]{images/Error_Analysis/XGB_IR_test_error_analysis_detailed.png}
    \end{minipage}
    \caption{Test-set error analysis for IR spectra across all four models.}
    \label{fig:error-analysis-ir}
\end{figure}

\subsubsection{IR}

Among all modalities, IR spectra yielded the most balanced error profiles. The
CNN achieved the strongest overall result with a macro-F1 score of 0.681 and a
mean per-label error rate of 35.3\%, closely followed by the MPS with a
macro-F1 of 0.671 and a mean error rate of 36.1\%. The TTN remained competitive
in terms of average error distribution, but its macro-F1 decreased to 0.564,
indicating less reliable performance on the more difficult and less frequent
labels. XGBoost showed the weakest IR performance, with both a lower macro-F1
score (0.537) and a clearly larger fraction of severely affected groups.

Despite the comparatively strong results for IR, several functional groups were
consistently difficult across models. In particular, thial and azo compound
showed very high missed-positive rates for all architectures, while hydrazone,
phosphine, sulfoxide, and acid anhydride also appeared repeatedly among the
most error-prone labels. This indicates that the remaining IR errors are not
distributed uniformly, but are concentrated in a small subset of chemically
challenging groups.

\begin{figure}[t]
    \centering
    \begin{minipage}[t]{0.49\linewidth}
        \centering
        \includegraphics[width=\linewidth]{images/Error_Analysis/CNN_H-NMR_test_error_analysis_detailed.png}
    \end{minipage}\hfill
    \begin{minipage}[t]{0.49\linewidth}
        \centering
        \includegraphics[width=\linewidth]{images/Error_Analysis/MPS_H-NMR_test_error_analysis_detailed.png}
    \end{minipage}

    \vspace{0.5em}

    \begin{minipage}[t]{0.49\linewidth}
        \centering
        \includegraphics[width=\linewidth]{images/Error_Analysis/TTN_H-NMR_test_error_analysis_detailed.png}
    \end{minipage}\hfill
    \begin{minipage}[t]{0.49\linewidth}
        \centering
        \includegraphics[width=\linewidth]{images/Error_Analysis/XGB_H_NMR_test_error_analysis_detailed.png}
    \end{minipage}
    \caption{Test-set error analysis for $^1$H-NMR spectra across all four models.}
    \label{fig:error-analysis-hnmr}
\end{figure}

\subsubsection{$^1$H-NMR}

The error analysis for $^1$H-NMR reveals a markedly more difficult prediction
task than IR. The CNN again produced the strongest overall result, combining the
highest macro-F1 score (0.493) with the lowest mean label-wise error rate
(59.7\%). The remaining models performed substantially worse in terms of class
balance. XGBoost reached a macro-F1 of 0.368, TTN 0.314, and MPS 0.256, while
all three showed broad error concentration across many labels. Notably, the MPS
displayed a mean error rate similar to the CNN, but its substantially lower
macro-F1 indicates that its remaining predictions were less evenly distributed
across functional groups.

Across all four models, sulfonic acid and thial were consistently among the
most difficult labels, frequently reaching error rates close to 100\%. Further
recurring failure cases included isothiocyanate, isocyanate, thiol, acid
anhydride, and hydrazone. Overall, the $^1$H-NMR results suggest that the
models can capture some global signal structure, but still struggle to separate
multiple chemically similar functional groups with sufficient robustness.

\begin{figure}[t]
    \centering
    \begin{minipage}[t]{0.49\linewidth}
        \centering
        \includegraphics[width=\linewidth]{images/Error_Analysis/CNN_C-NMR_test_error_analysis_detailed.png}
    \end{minipage}\hfill
    \begin{minipage}[t]{0.49\linewidth}
        \centering
        \includegraphics[width=\linewidth]{images/Error_Analysis/MPS_C-NMR_test_error_analysis_detailed.png}
    \end{minipage}

    \vspace{0.5em}

    \begin{minipage}[t]{0.49\linewidth}
        \centering
        \includegraphics[width=\linewidth]{images/Error_Analysis/TTN_C-NMR_test_error_analysis_detailed.png}
    \end{minipage}\hfill
    \begin{minipage}[t]{0.49\linewidth}
        \centering
        \includegraphics[width=\linewidth]{images/Error_Analysis/XGB_C_NMR_test_error_analysis_detailed.png}
    \end{minipage}
    \caption{Test-set error analysis for $^{13}$C-NMR spectra across all four models.}
    \label{fig:error-analysis-cnmr}
\end{figure}

\subsubsection{$^{13}$C-NMR}

The $^{13}$C-NMR setting exhibits a mixed picture and highlights the importance
of complementing global scores with label-wise error analysis. XGBoost achieved
the best macro-F1 score in this modality (0.413), while MPS showed the lowest
mean per-label error rate (58.1\%). TTN followed closely behind with a mean
error of 63.3\%, whereas the CNN performed weakest overall, with a macro-F1 of
only 0.150 and a mean error rate of 82.2\%. This divergence between macro-F1
and average per-label error shows that some models are better at maintaining
balanced predictions across labels, while others produce lower average error
rates but remain inconsistent for specific functional groups.

The most difficult $^{13}$C-NMR labels were again highly consistent across
models. Sulfonic acid and thial were universally problematic, and high error
rates were also observed for isothiocyanate, azo compound, phosphine, thiol,
and isocyanate. Compared with IR, the error distributions were shifted more
strongly toward high values, indicating that $^{13}$C-NMR alone provides less
discriminative information for a large fraction of the target groups.

\begin{figure}[t]
    \centering
    \begin{minipage}[t]{0.49\linewidth}
        \centering
        \includegraphics[width=\linewidth]{images/Error_Analysis/CNN_MSMS+_test_error_analysis_detailed.png}
    \end{minipage}\hfill
    \begin{minipage}[t]{0.49\linewidth}
        \centering
        \includegraphics[width=\linewidth]{images/Error_Analysis/MPS_MSMS+_test_error_analysis_detailed.png}
    \end{minipage}

    \vspace{0.5em}

    \begin{minipage}[t]{0.49\linewidth}
        \centering
        \includegraphics[width=\linewidth]{images/Error_Analysis/TTN_MSMS+_test_error_analysis_detailed.png}
    \end{minipage}\hfill
    \begin{minipage}[t]{0.49\linewidth}
        \centering
        \includegraphics[width=\linewidth]{images/Error_Analysis/XGB_MSMS+_test_error_analysis_detailed.png}
    \end{minipage}
    \caption{Test-set error analysis for positive-ion MS/MS spectra across all four models.}
    \label{fig:error-analysis-msms-pos}
\end{figure}

\subsubsection{MS/MS (+)}

For positive-ion MS/MS spectra, the error analysis reveals a clear separation
between MPS and the remaining three models. MPS achieved the best macro-F1
score (0.175) and by far the lowest mean per-label error rate (44.8\%), whereas
CNN, TTN, and XGBoost all exhibited mean error rates above 85\% and extremely
high medians close to 100\%. In practical terms, this means that for these
three models the majority of functional groups were missed almost completely on
the test set, even though their micro-F1 scores remained superficially higher
than the macro-F1 values would suggest.

The MPS error profile for MS/MS (+) was still far from ideal, but it was much
more graded and therefore more interpretable than those of the other models.
Its most difficult labels included sulfonic acid, hydrazone, enol, alcohol,
alkene, and phenol. By contrast, CNN, TTN, and XGBoost assigned near-maximal
error rates to a large portion of the label space, indicating that these models
failed to extract sufficiently transferable representations from the positive-ion
fragmentation patterns.

\begin{figure}[t]
    \centering
    \begin{minipage}[t]{0.49\linewidth}
        \centering
        \includegraphics[width=\linewidth]{images/Error_Analysis/CNN_MSMS-_test_error_analysis_detailed.png}
    \end{minipage}\hfill
    \begin{minipage}[t]{0.49\linewidth}
        \centering
        \includegraphics[width=\linewidth]{images/Error_Analysis/MPS_MSMS-_test_error_analysis_detailed.png}
    \end{minipage}

    \vspace{0.5em}

    \begin{minipage}[t]{0.49\linewidth}
        \centering
        \includegraphics[width=\linewidth]{images/Error_Analysis/TTN_MSMS-_test_error_analysis_detailed.png}
    \end{minipage}\hfill
    \begin{minipage}[t]{0.49\linewidth}
        \centering
        \includegraphics[width=\linewidth]{images/Error_Analysis/XGB_MSMS-_test_error_analysis_detailed.png}
    \end{minipage}
    \caption{Test-set error analysis for negative-ion MS/MS spectra across all four models.}
    \label{fig:error-analysis-msms-neg}
\end{figure}

\subsubsection{MS/MS (-)}

The negative-ion MS/MS results closely mirror the behavior observed for MS/MS
(+). MPS again achieved the most favorable error distribution, with the best
macro-F1 score (0.171) and the lowest mean per-label error rate (43.2\%).
CNN, TTN, and XGBoost again showed a near-collapse of label-wise performance,
with mean error rates between 86.1\% and 89.4\% and medians close to 100\%.
This indicates that the weak aggregate performance of these models on MS/MS is
not caused by a few isolated outlier labels, but by a broad and systematic
failure across most of the functional group space.

Within the MPS model, the main sources of error for MS/MS (-) were sulfonic
acid, hydrazone, alcohol, ketone, sulfide, phenol, imine, and ester. While this
still reflects substantial difficulty, the distribution was considerably less
extreme than for the competing approaches. Taken together, the MS/MS analyses
show that MPS is the only model in this comparison that retained a meaningful
degree of discrimination for both ionization modes.

Overall, the test-set error analysis confirms three central observations. First,
IR is the most informative and most stable modality, with CNN and MPS producing
the most balanced results. Second, $^1$H-NMR and $^{13}$C-NMR are substantially
more challenging, with performance strongly depending on the model and with
recurrent failure cases for rare and chemically specific labels. Third, the
MS/MS modalities are by far the hardest setting in this benchmark, and only MPS
maintains a label-wise error profile that remains meaningfully interpretable.
Across nearly all modalities, labels such as sulfonic acid, thial, azo
compound, isothiocyanate, hydrazone, and phosphine repeatedly appeared among the
most difficult groups, suggesting that these functional groups represent a
persistent source of ambiguity across both classical and tensor-network-based
approaches.
